\section{Literature Review and Related Work}
\label{sec:related}

\subsection{Fashion-MNIST}
Xiao, Rasul and Vollgraf introduce Fashion-MNIST, a set of 70{,}000 grayscale
$28\times 28$ images of fashion products in 10 categories, with 7{,}000 images per category
and an official split of 60{,}000 training and 10{,}000 test images~\cite{xiao2017fashion}.
They present it as a direct drop-in replacement for MNIST: it shares the image size, data
format and split structure, so existing MNIST code needs essentially only a different download
location. Their motivation is that MNIST is popular because it is small, simply encoded and
supported by common libraries, but it is an easy task; Fashion-MNIST aims to keep those
advantages while being more challenging~\cite{xiao2017fashion}.

The images are front-view thumbnails of 70{,}000 unique products from an online fashion shop,
converted to $28\times 28$ grayscale by a seven-step pipeline. Labels are product silhouette
codes assigned manually by in-house experts and reviewed by a separate team. The files use the
MNIST format and are stored sorted by label, so shuffling is left to the user
\cite{xiao2017fashion}; our train/validation split is therefore stratified and seeded. The
benchmark in~\cite{xiao2017fashion} covers scikit-learn classifiers and contains no
convolutional network.

\subsection{Learned versus Hand-Designed Features}
LeCun et al.\ review methods for handwritten character recognition and compare them on a
standard handwritten digit task~\cite{lecun1998gradient}. Their abstract argues that
multilayer networks trained by back-propagation can learn complex decision surfaces on
high-dimensional patterns with minimal preprocessing, and reports that convolutional networks,
which are designed for the variability of 2D shapes, outperform the other techniques compared
in that study~\cite{lecun1998gradient}. This motivates the learned-feature side of our
comparison.

\subsection{Relation to This Work}
We compare a small CNN, which learns spatial features, against four classical models
(logistic regression, $k$-NN, random forest, RBF-SVM) on two fixed representations: flattened
pixels and HOG features. Neither cited paper ran this experiment: \cite{xiao2017fashion}
does not evaluate HOG or a CNN, and \cite{lecun1998gradient} studies handwritten digits and
documents rather than clothing images. We therefore draw no conclusion from their reported
results; our findings are the measurements in Section~\ref{sec:results}.
